% ##################################################
% Gemeinsames Setup fuer die Probeklausuren
% (Java-Listings, Hinweis-/Beispiel-Boxen, Struktogramm-Stile in TikZ)
% wird nach \include{../for_latex/preamble} per \input eingebunden
% ##################################################

% --- Java Code-Highlighting ---
\definecolor{codegray}{rgb}{0.5,0.5,0.5}
\definecolor{codepurple}{rgb}{0.58,0,0.82}
\definecolor{backcolour}{rgb}{0.95,0.95,0.92}
\definecolor{keywordcolor}{rgb}{0.8,0.47,0.13}

\lstdefinestyle{javastyle}{
	backgroundcolor=\color{backcolour},
	commentstyle=\color{codegray},
	keywordstyle=\color{keywordcolor}\bfseries,
	numberstyle=\tiny\color{codegray},
	stringstyle=\color{codepurple},
	basicstyle=\ttfamily\footnotesize,
	breakatwhitespace=false,
	breaklines=true,
	captionpos=b,
	keepspaces=true,
	numbers=left,
	numbersep=5pt,
	showspaces=false,
	showstringspaces=false,
	showtabs=false,
	tabsize=2,
	language=Java,
	frame=single,
	rulecolor=\color{black!30}
}

\lstset{style=javastyle}

% --- Hinweis-/Beispiel-/Formel-Umgebungen (schlicht, ohne farbige Box) ---
\newenvironment{hinweis}[1][Hinweis]%
	{\par\smallskip\noindent\textbf{#1.}\enspace\ignorespaces}%
	{\par\smallskip}

\newenvironment{beispiel}[1][Beispiel]%
	{\par\smallskip\noindent\textbf{#1:}\par\nobreak\vspace{2pt}}%
	{\par\smallskip}

\newenvironment{formeln}[1][Gegebene Formeln]%
	{\par\smallskip\noindent\textbf{#1:}\par\nobreak\vspace{2pt}}%
	{\par\smallskip}

% ##################################################
% Struktogramm (Nassi-Shneiderman) in reinem TikZ
% ##################################################
\newlength{\nswidth}\setlength{\nswidth}{11.5cm}   % Gesamtbreite
\newlength{\nsind}\setlength{\nsind}{0.55cm}        % Einrueckung je Schleifenebene
% vorberechnete Versatzbreiten fuer Verzweigungen (in Koordinaten nutzbar)
\newlength{\nshalf}\setlength{\nshalf}{\dimexpr\nswidth/2\relax}            % halbe Gesamtbreite
\newlength{\nshalfi}\setlength{\nshalfi}{\dimexpr(\nswidth-\nsind)/2\relax}  % halbe Breite, 1x eingerueckt
\newlength{\nshalfii}\setlength{\nshalfii}{\dimexpr(\nswidth-2\nsind)/2\relax} % halbe Breite, 2x eingerueckt

\tikzset{
	% Basis fuer Anweisungs-/Verzweigungsbloecke
	nb/.style={draw, line width=0.4pt, anchor=north west, inner sep=4pt,
		align=left, font=\footnotesize, fill=white, minimum height=0.8cm},
	% volle Breite
	nbfull/.style={nb, minimum width=\nswidth,
		text width=\dimexpr\nswidth-8pt\relax},
	% eine Ebene eingerueckt (Schleifenrumpf)
	nbi/.style={nb, minimum width=\dimexpr\nswidth-\nsind\relax,
		text width=\dimexpr\nswidth-\nsind-8pt\relax},
	% zwei Ebenen eingerueckt (Schleife in Schleife)
	nbii/.style={nb, minimum width=\dimexpr\nswidth-2\nsind\relax,
		text width=\dimexpr\nswidth-2\nsind-8pt\relax},
	% Verzweigung: halbe Breiten
	nbh/.style={nb, minimum width=\dimexpr\nswidth/2\relax,
		text width=\dimexpr\nswidth/2-8pt\relax},
	nbih/.style={nb, minimum width=\dimexpr(\nswidth-\nsind)/2\relax,
		text width=\dimexpr(\nswidth-\nsind)/2-8pt\relax},
	% Verzweigung halbe Breite, zwei Ebenen eingerueckt
	nbiih/.style={nb, minimum width=\nshalfii,
		text width=\dimexpr\nshalfii-8pt\relax},
	% Kopf einer Verzweigung (das Dreieck), inner sep 0
	nc/.style={draw, line width=0.4pt, anchor=north west, inner sep=0pt,
		minimum height=1.0cm, fill=white},
	ncfull/.style={nc, minimum width=\nswidth},
	nci/.style={nc, minimum width=\dimexpr\nswidth-\nsind\relax},
	ncii/.style={nc, minimum width=\dimexpr\nswidth-2\nsind\relax},
}

% Zeichnet das Dreieck + Beschriftungen in einen bereits gesetzten nc-Knoten.
% #1 Knotenname  #2 Bedingung  #3 ja-Label  #4 nein-Label
\newcommand{\nscond}[4]{%
	\draw[line width=0.4pt] (#1.north west) -- (#1.south);
	\draw[line width=0.4pt] (#1.north east) -- (#1.south);
	\node[font=\scriptsize, anchor=north, inner sep=1pt] at ($(#1.north)+(0,-1.5pt)$) {#2};
	\node[font=\scriptsize, anchor=south west, inner sep=1pt] at ($(#1.south west)+(2.5pt,1.5pt)$) {#3};
	\node[font=\scriptsize, anchor=south east, inner sep=1pt] at ($(#1.south east)+(-2.5pt,1.5pt)$) {#4};
}

% Zeichnet den L-Rahmen einer Schleife (linker + unterer Rand der Einrueckspalte).
% #1 Kopfknoten der Schleife  #2 letzter Rumpfknoten
\newcommand{\nsframe}[2]{%
	\draw[line width=0.4pt] (#1.south west) |- (#2.south west);
}
